\documentclass[12pt,a4paper]{article}

\usepackage[margin=2.5cm]{geometry}

\title{VVSI Project}
\author{}
\date{}

\begin{document}

\maketitle

\begin{abstract}
% Maximum one page. Write this section last.
\end{abstract}

\section{Introduction}
% Write this section last.

\subsection{Context}
One of the biggest challenges in finance is to be able to accurately predict the prices of a asset, it being a stock paper, options, crypto currencies, commodities, any type of product offered in the market. The reason is simple: if you know the price of the next day, you are able to fully optimize your profit. The base idea on predicting the price is simple: we want to know the price in the future (one day, five days, ...) based on the historical data of that asset, which is characterized as a time-series. 

Time-series analysis is a complex type of investigation, since it involves many aspects such as stationarity, ergodicity, roughness, seasonality etc. Statistics have been tackling the problem of historical forecasting by analyzing these characteristics. However, with the rapid evolution of ML (Machine Learning) methods in the past 20 years, especially, DNN (Deep Neural Networks), with LSTMs and LLMs, there has been a risen on leveraging these models for performing such analysis, while delivering results that are better or even equal to statistical methods. 

\subsection{Motivations}
In the context of ML, the data is the key for a good model. However, the data is not explored as it's in the Statistics field, where the main focus is understanding properties of the data and devising a model dependent on those for forecasting the future and focused on one object of study (local model). The models in ML are built to derive the characteristics from the data and consequently perform the forecasting after a training process, while still maintaing the capability of being expanded for more than one subject (global model). 

In this study, we are interested in understading how well the base LSTM predicts the overnight return of stock papers, while being trained with data composed of multiple tickets. 

\subsection{Contributions}
In this project, we have two main contributions: 
\begin{itemize}
    \item A reproducible pipeline for collecting and preparing 20 years of historical data across 1,834 equities.
    \item A global LSTM model that predicts next-day overnight returns across many companies.

\end{itemize}

\subsection{Related Work}
Legacy econometric models as GARCH \cite{bollerslev1986} and SARIMA \cite{box2015} have been extensive used for the purpose of predicting equities prices due to their instrinsic capacity of accounting for data-specific properties such as seasonality, stationarity, mean-reversal, etc. However, they lack the necessary tools for explaining non-linear behaviour.

LSTMs \cite{vennerod2021} have been widely used as an substitute of RNNs due to its architecture amenizing the gradient-vanishing problem, while still maintaing an accurate cronological representation. Nevertheless, great majority of its implementations are with local models .

\subsection{Outline}
 \begin{enumerate}
      \item \textbf{Section 2 - Background} introduces the financial and
      machine-learning concepts required to understand the project.

      \item \textbf{Section 3 - Methods} describes the dataset preparation,
      forecasting problem, LSTM architecture, data splits, and evaluation
      methodology.

      \item \textbf{Section 4 - Implementation} presents the implementation
      of the data pipeline, model-training process, checkpoint selection, and
      experiment tracking.

      \item \textbf{Section 5 - Experimental Results} describes the
      experimental setting and compares the results obtained with input windows
      of 10, 30, and 100 trading days.

      \item \textbf{Section 6 - Conclusions} summarizes the findings,
      discusses the limitations of the project, and proposes possible future
      developments.
  \end{enumerate}

\section{Background}
% Present all background knowledge required to understand the work.

\section{Methods}
% Describe the system design and methodology. Do not include source code here.

\section{Implementation}
% Describe how the system was implemented.
% Use pseudocode and/or informal or UML-style diagrams where useful.

\section{Experimental Results}
% Describe the experimental results. This is the most important section.

\subsection{Goals}
% Describe the objectives of the experimental activity, typically:
% - demonstrate correctness;
% - evaluate computational performance, such as CPU/GPU time and RAM usage;
% - evaluate effectiveness, such as expected gain and its standard deviation.

\subsection{Setting}
% Describe the experimental setting: hardware, software, datasets, splits,
% hyperparameters, reproducibility controls, and evaluation metrics.

\subsection{Case Studies}
% Describe the case studies used in the experiments.

\subsection{Correctness}
% Describe experiments that demonstrate the correctness of the implementation.

\subsection{Computational Evaluation}
% Evaluate computational performance, such as training time and CPU, GPU,
% RAM, and VRAM usage.

\subsection{Technical Evaluation}
% Evaluate the effectiveness of the proposed approach.
% Describe the experiments and summarize the results using tables and graphs.

\section{Conclusions}
% Summarize the work and outline possible future developments.

\begin{thebibliography}{99}
\bibitem{bollerslev1986}
T. Bollerslev,
``Generalized autoregressive conditional heteroskedasticity,''
\textit{Journal of Econometrics}, vol. 31, no. 3, pp. 307--327, 1986.

\bibitem{box2015}
G. E. P. Box, G. M. Jenkins, G. C. Reinsel, and G. M. Ljung,
\textit{Time Series Analysis: Forecasting and Control}, 5th ed.
Hoboken, NJ: John Wiley \& Sons, 2015.

\bibitem{vennerod2021}
C. B. Venner{\o}d, A. Kj{\ae}rran, and E. S. Bugge,
``Long short-term memory RNN,'' arXiv:2105.06756, 2021.
\end{thebibliography}

\end{document}
